Una càrrega puntual $Q_1 = +1,00\times 10^{-8}\ \mathrm{C}$ està situada a l'origen de coordenades. Una altra càrrega puntual $Q_2 = -2,00\times 10^{-8}\ \mathrm{C}$ està situada en el semieix $Y$ positiu, a 3,00~m de l'origen. Calculeu:

\begin{apartat}{1}
El camp i el potencial electrostàtic en un punt A situat en el semieix $X$ positiu, a 4,00~m de l'origen. Dibuixeu un esquema de tots els camps elèctrics que intervenen en el problema.
\end{apartat}

\begin{apartat}{1}
El treball fet pel camp elèctric en traslladar una càrrega puntual d'1,00~C des del punt A a un punt B de coordenades $(4,00,\ 3,00)$~m.
\end{apartat}

\blocetiquetat{Dada:}{$k = 1/4\pi\varepsilon_0 = 8,99\times 10^{9}\ \mathrm{N\,m^2\,C^{-2}}$}
